\documentclass[12pt, a4paper]{article}

\usepackage[utf8]{inputenc}
\usepackage[T2A]{fontenc}
\usepackage[russian]{babel}
\usepackage{amsmath, amssymb, amsthm}
\usepackage{mathtools}
\usepackage{geometry}
\usepackage{xcolor}
\usepackage{tcolorbox}
\usepackage{booktabs}
\usepackage{array}
\usepackage{enumitem}
\usepackage{hyperref}
\usepackage{fancyhdr}
\usepackage{parskip}

\tcbuselibrary{skins, breakable}

\geometry{left=2.5cm, right=2.5cm, top=2.5cm, bottom=2.5cm}

% ===== ЦВЕТА =====
\definecolor{defblue}{RGB}{13, 71, 161}
\definecolor{defbluebg}{RGB}{227, 242, 253}
\definecolor{thmgreen}{RGB}{27, 94, 32}
\definecolor{thmgreenbg}{RGB}{232, 245, 233}
\definecolor{propviolet}{RGB}{74, 20, 140}
\definecolor{propvioletbg}{RGB}{243, 229, 245}
\definecolor{remamber}{RGB}{180, 60, 0}
\definecolor{remamberbg}{RGB}{255, 243, 224}
\definecolor{exgraybg}{RGB}{245, 245, 245}
\definecolor{exgray}{RGB}{60, 60, 60}
\definecolor{corollteal}{RGB}{0, 77, 64}
\definecolor{corolltbg}{RGB}{224, 242, 241}
\definecolor{proofbg}{RGB}{250, 250, 250}
\definecolor{qcolor}{RGB}{130, 0, 30}

% ===== БЛОКИ =====
\tcbset{
  mybase/.style={
    breakable, enhanced,
    fonttitle=\bfseries\small,
    left=8pt, right=8pt, top=5pt, bottom=5pt,
    arc=4pt, boxrule=1pt,
    parbox=false
  }
}

\newtcolorbox{defbox}[1]{
  mybase,
  colback=defbluebg, colframe=defblue,
  colbacktitle=defblue, coltitle=white,
  title={\small Определение\ifx&#1&\else{}: {#1}\fi}
}
\newtcolorbox{thmbox}[1]{
  mybase,
  colback=thmgreenbg, colframe=thmgreen,
  colbacktitle=thmgreen, coltitle=white,
  title={\small Теорема\ifx&#1&\else{}: {#1}\fi}
}
\newtcolorbox{propbox}[1]{
  mybase,
  colback=propvioletbg, colframe=propviolet,
  colbacktitle=propviolet, coltitle=white,
  title={\small Утверждение\ifx&#1&\else{}: {#1}\fi}
}
\newtcolorbox{rembox}[1]{
  mybase,
  colback=remamberbg, colframe=remamber,
  colbacktitle=remamber, coltitle=white,
  title={\small Замечание\ifx&#1&\else{}: {#1}\fi}
}
\newtcolorbox{exbox}[1]{
  mybase,
  colback=exgraybg, colframe=exgray,
  colbacktitle=exgray, coltitle=white,
  title={\small Пример\ifx&#1&\else{}: {#1}\fi}
}
\newtcolorbox{corbox}[1]{
  mybase,
  colback=corolltbg, colframe=corollteal,
  colbacktitle=corollteal, coltitle=white,
  title={\small Следствие\ifx&#1&\else{}: {#1}\fi}
}
\newtcolorbox{proofbox}{
  mybase,
  colback=proofbg, colframe=gray!50,
  colbacktitle=gray!50, coltitle=black,
  title={\small Доказательство}
}

% ===== ЗАГОЛОВКИ =====
\newcommand{\question}[2]{%
  \newpage%
  \noindent\colorbox{qcolor}{\parbox{\dimexpr\linewidth-2\fboxsep\relax}{%
    \color{white}\large\bfseries\strut Вопрос #1. #2%
  }}%
  \vspace{6pt}%
}
\newcommand{\qsub}[1]{%
  \vspace{6pt}%
  {\bfseries\color{defblue} #1}%
  \par\vspace{2pt}%
}

% ===== КОЛОНТИТУЛЫ =====
\pagestyle{fancy}
\fancyhf{}
\fancyhead[L]{\small\color{gray} Статистика для МО --- Блок 1 (Вопросы 1--8)}
\fancyhead[R]{\small\color{gray} ИТМО, 2026}
\fancyfoot[C]{\small\thepage}
\renewcommand{\headrulewidth}{0.4pt}
\setlength{\headheight}{14pt}

% ===== МАТ. МАКРОСЫ =====
\newcommand{\cov}{\operatorname{cov}}
\newcommand{\Var}{\operatorname{Var}}
\newcommand{\E}{\mathbb{E}}
\newcommand{\sign}{\operatorname{sign}}
\newcommand{\rk}{\operatorname{rk}}
\newcommand{\supp}{\operatorname{supp}}

\hypersetup{colorlinks=true, linkcolor=defblue, urlcolor=defblue, hidelinks}
\setlist[enumerate]{leftmargin=*, topsep=2pt, itemsep=1pt}
\setlist[itemize]{leftmargin=*, topsep=2pt, itemsep=1pt}

\begin{document}

%%%%%%%%%%%%%%%%%%%%%%%%%%%%%%%%%%%%%%%%%%%%%%%%%%%%%
\question{1}{Гауссовские векторы и их свойства}
%%%%%%%%%%%%%%%%%%%%%%%%%%%%%%%%%%%%%%%%%%%%%%%%%%%%%

\qsub{Стандартный гауссовский вектор}

\begin{defbox}{Стандартный гауссовский вектор}
Вектор $X = (X_1, \ldots, X_n)^T$ называется \textbf{стандартным гауссовским}, если
\[
  X_i \sim \mathcal{N}(0,1), \quad X_1,\ldots,X_n \text{ --- независимы.}
\]
Обозначение: $X \sim \mathcal{N}(0, I)$.

Совместная плотность:
\[
  p(x_1,\ldots,x_n)
  = \prod_{i=1}^{n}\frac{1}{\sqrt{2\pi}}\,e^{-\frac{x_i^2}{2}}
  = \frac{1}{(\sqrt{2\pi})^n}\,e^{-\frac{1}{2}x^Tx}.
\]
\end{defbox}

\qsub{Матрица ковариаций}

\begin{defbox}{Матрица ковариаций}
Для случайного вектора $X$:
\[
  \Var X = \bigl(\cov(X_i,X_j)\bigr)_{i,j} = \E\bigl[(X-\E X)(X-\E X)^T\bigr].
\]
\end{defbox}

\begin{propbox}{Преобразование матрицы ковариаций}
Если $Y = BX$, то $\Var(BX) = B\,\Var X\,B^T$.
\end{propbox}

\begin{proofbox}
\[
  \Var(BX) = \E[(BX - \E BX)(BX-\E BX)^T]
  = \E[B(X-\E X)(X-\E X)^T B^T]
  = B\,\Var X\,B^T. \quad\square
\]
\end{proofbox}

\begin{thmbox}{Независимость $\Leftrightarrow$ некоррелированность для гауссовских}
Если $(X,Y)\sim\mathcal{N}(\mu,\Sigma)$, то
\[
  X \text{ и } Y \text{ независимы} \;\iff\; \cov(X,Y) = 0.
\]
\end{thmbox}

\begin{proofbox}
Б.о.о. $\E X = \E Y = 0$. Пусть $\cov(X,Y)=0$, тогда
$\Sigma = \begin{pmatrix}\sigma_X^2 & 0\\0 & \sigma_Y^2\end{pmatrix}$.
\begin{align*}
  p(x,y)
  &= \frac{1}{2\pi\sigma_X\sigma_Y}
     \exp\!\left(-\tfrac{1}{2}\!\left(\tfrac{x^2}{\sigma_X^2}+\tfrac{y^2}{\sigma_Y^2}\right)\right)
   = \underbrace{\frac{1}{\sqrt{2\pi}\,\sigma_X}e^{-\frac{x^2}{2\sigma_X^2}}}_{p_X(x)}
   \cdot
     \underbrace{\frac{1}{\sqrt{2\pi}\,\sigma_Y}e^{-\frac{y^2}{2\sigma_Y^2}}}_{p_Y(y)}.
\end{align*}
Совместная плотность разложилась в произведение маргинальных $\Rightarrow$ независимость. $\square$
\end{proofbox}

\begin{rembox}{}
Этот факт специфичен для гауссовского случая. В общем случае некоррелированность
\textbf{не} влечёт независимость.
\end{rembox}

%%%%%%%%%%%%%%%%%%%%%%%%%%%%%%%%%%%%%%%%%%%%%%%%%%%%%
\question{2}{Независимость и некоррелированность}
%%%%%%%%%%%%%%%%%%%%%%%%%%%%%%%%%%%%%%%%%%%%%%%%%%%%%

\begin{defbox}{Независимость}
$X$ и $Y$ \textbf{независимы}, если:
\[
  F_{X,Y}(x,y) = F_X(x)\cdot F_Y(y).
\]
\end{defbox}

\begin{defbox}{Некоррелированность}
$X$ и $Y$ \textbf{некоррелированы}, если
\[
  \rho(X,Y) = \frac{\cov(X,Y)}{\sqrt{\Var X \cdot \Var Y}} = 0.
\]
\end{defbox}

\begin{exbox}{Некоррелированные, но зависимые с.в.}
$(X,Y)$ равномерно на 4 точках $(1,0),(-1,0),(0,1),(0,-1)$ с $p=1/4$.

$\E[XY] = 0$, $\E X = \E Y = 0 \Rightarrow \cov = 0$.

Но $P(X=1,Y=1) = 0 \neq 1/16 = P(X=1)\cdot P(Y=1)$ --- зависимы.
\end{exbox}

\end{document}
